\clearpage
\section{Learning token representations}\label{sec:embeddings}
\subsection*{From context counts to association}
Words that occur in similar contexts can acquire related representations.
The choice of ``context'' matters: neighboring words, a document, and a
syntactic relation define different statistics. In the lecture's constructed
word--context table, \texttt{tea} and \texttt{coffee} share \texttt{drink}
and \texttt{hot}, while \texttt{car} occurs with \texttt{drive}.

Let $X_{wc}$ be the count for word $w$ and context $c$, and
$Z=\sum_{w,c}X_{wc}$. Pointwise mutual information compares the observed
joint frequency with the frequency expected under independence:
\[
 \operatorname{PMI}(w,c)=\log_2\frac{X_{wc}Z}{X_{w\cdot}X_{\cdot c}},\qquad
 \operatorname{PPMI}(w,c)=\max(0,\operatorname{PMI}(w,c)).
\]
Here the row and column totals are positive. An unobserved pair with positive
marginals has PMI $-\infty$ and is assigned PPMI zero.

\fig{counts-ppmi}{The L03 invented count table and its PPMI transform. Counts
sum to 35. The transformed values measure association in bits, not probability.}

For \texttt{tea}/\texttt{drink}, the joint count is eight, the row total is
14, and the column total is 15. Therefore
\[
 \operatorname{PMI}(\texttt{tea},\texttt{drink})
 =\log_2\frac{8\cdot35}{14\cdot15}
 =\log_2(4/3)\approx0.4150\text{ bits}.
\]
The zero entries after PPMI conflate unobserved and nonpositively associated
pairs. This choice is part of the representation, not a lossless rewrite of counts.

\subsection*{Compression and alternative objectives}
For a matrix $M$, truncated SVD forms $M_r=U_r\Sigma_rV_r^\top$.
It minimizes squared reconstruction error among matrices of rank at most
$r$. A word table can use $U_r\Sigma_r$, producing a dense vector for each
row; the scaling convention should be stated because it affects geometry.

GloVe instead fits weighted log counts using vectors and biases:
\[
 J=\sum_{X_{wc}>0}f(X_{wc})
 \bigl(u_w^\top v_c+b_w+\widetilde b_c-\log X_{wc}\bigr)^2.
\]
The weight function reduces the influence of selected count ranges.
This objective differs from applying SVD to PPMI. Both produce vectors,
but what they preserve is determined by their input statistics and loss.

\selfcheck{Why can a frequent word have low PMI with a frequent context?
The independence baseline is already large. See \readingref{R07} for GloVe's objective.}

\nextpage{Predictive embeddings and one gradient by hand}
CBOW predicts a center word from its surrounding context; skip-gram predicts
surrounding contexts from a center word. These are representation-learning
tasks. A symmetric context window may contain future words, so it should
not be mistaken for the prefix-only conditioning of a causal language model.

\fig{cbow-skipgram}{Prediction direction changes the training examples.
Word2vec uses distinct input and output roles even when both index the same vocabulary.}

Negative sampling replaces a vocabulary-wide prediction objective with
observed-versus-sampled pair discrimination. For one positive context $c_+$
and one sampled negative context $c_-$, write
\[
 \ell=-\log\sigma(w^\top c_+)-\log\sigma(-w^\top c_-),\qquad
 \sigma(a)=\frac1{1+e^{-a}}.
\]
The sigmoid values do not sum to one across vocabulary items. This loss
is therefore not a normalized next-token negative log likelihood.

\fig{negative-gradient}{The positive score is one and the negative score is
$0.5$. Gradient descent increases the positive score when the center vector is held fixed.}

Using $w=(1,0.5)$, $c_+=(0.5,1)$, and $c_-=(1,-1)$ gives
$\sigma(1)\approx0.7311$, $\sigma(0.5)\approx0.6225$, and
$\ell\approx1.2873$. Differentiating the positive term gives
\[
 \nabla_{c_+}\ell=(\sigma(w^\top c_+)-1)w
 \approx(-0.2689,-0.1345).
\]
At step size $\eta=0.5$, the positive context becomes approximately
$(0.6345,1.0672)$. With $w$ fixed, its score rises from 1 to 1.1681.
The negative-context gradient is $\sigma(w^\top c_-)w$; descent reduces
that score. The center vector receives contributions from both pairs.

\takeaway{Specify the prediction task and objective before interpreting
the vectors. The sampling distribution also influences what the model learns.}
\selfcheck{At zero scores, each binary term contributes $\log2$. What is the
initial loss with one positive and five negatives? It is $6\log2\approx4.1589$.}
\textbf{Reading connection.} \readingref{R05} introduces CBOW/skip-gram;
\readingref{R06}, Section 2.2, gives the negative-sampling objective.

\nextpage{Static vectors, contextual states, and subword composition}
An embedding table is a learned collection of rows. Looking up the ID for
\texttt{bank} returns the same row every time that ID occurs. Contextual
computation can transform that row together with other inputs into a
different state for ``river bank'' and ``bank loan.'' Learning a good table
does not by itself supply that context computation.

\fig{representation-types}{A token row, a contextual state, and a retrieval
vector have different inputs and uses. A retrieval representation may summarize
an entire passage and be trained for matching.}

Cosine similarity compares the directions of two nonzero vectors:
\[
 \operatorname{cos}(u,v)=\frac{u^\top v}{\|u\|_2\|v\|_2}.
\]
Nearest neighbors and vector analogies are useful inspections of geometry.
They depend on the corpus, objective, normalization, and chosen examples.
A plausible neighbor list does not establish language-model likelihood,
factual accuracy, or performance on a downstream task.

\fig{subword-composition}{An illustrative fastText-style composition sums
overlapping character $n$-gram vectors. BPE segmentation supplies a sequence
of token IDs. The two operations use subwords in different ways.}

In a character $n$-gram representation, a word can be represented by a
word-specific vector plus vectors for its subword features. Overlapping
features can share statistical strength between related forms, and known
features can help represent a previously unseen word. The exact boundary
symbols, $n$-gram lengths, hashing, and normalization belong to the chosen
implementation. The figure only illustrates the composition principle.

In a BPE-based language model, the tokenizer instead emits separate subword
tokens. Each token has a lookup row; the model then processes their ordered
sequence. Character-feature composition and BPE can address related coverage
problems, but they define different model inputs.

\textbf{What carries forward.} Modern neural LMs still begin with lookup
tables and learn output scores. The additional question is how the model
forms a state from its available context. We next make the lookup, target
alignment, objective, and parameter updates explicit.

\selfcheck{If two occurrences have the same token ID, when must their vectors
be identical? At a deterministic table lookup. After a context-dependent
network, their states can differ.}
